\begin{table*}[t]\centering\footnotesize
\caption{Detail recovery by land-cover class on SPOT$\rightarrow$aerial, over 300 held-out tiles. \emph{Target} is the mean gradient magnitude of the aerial truth over pixels of that class---how much high-frequency content exists to recover. The model columns give the fraction of it recovered. Classes are ordered by share of test pixels.}
\label{tab:perclass}
\setlength{\tabcolsep}{4pt}
\begin{tabular}{@{}lrrrrr@{}}\toprule
Class & Share & Target & Bicubic & EDSR & +\,nDSM\\\midrule
Agricultural land & 43.7\% & 0.0260 & 3.3\% & 15.3\% & \textbf{18.2\%}\\
Deciduous & 23.8\% & 0.0271 & 4.9\% & 13.1\% & \textbf{13.6\%}\\
Herbaceous vegetation & 9.7\% & 0.0210 & 8.1\% & 21.3\% & \textbf{23.1\%}\\
Pervious surface & 5.7\% & 0.0319 & 10.3\% & 20.3\% & \textbf{22.8\%}\\
Water & 4.3\% & 0.0070 & 12.8\% & 53.7\% & \textbf{81.1\%}\\
Greenhouse & 3.9\% & 0.0280 & 14.2\% & 20.3\% & \textbf{22.1\%}\\
Plowed land & 2.4\% & 0.0253 & 4.3\% & 15.5\% & \textbf{13.0\%}\\
Impervious surface & 2.2\% & 0.0320 & 11.1\% & 21.5\% & \textbf{24.0\%}\\
\bottomrule\end{tabular}\end{table*}